\documentclass[11pt]{article}

\usepackage{amsmath,amssymb,amsthm,mathtools}
\usepackage[margin=1in]{geometry}
\usepackage{hyperref,microtype}

\newtheorem{theorem}{Theorem}
\newtheorem{lemma}[theorem]{Lemma}
\newtheorem{proposition}[theorem]{Proposition}
\newtheorem{corollary}[theorem]{Corollary}
\theoremstyle{definition}
\newtheorem{definition}[theorem]{Definition}
\newtheorem{remark}[theorem]{Remark}
\newtheorem{construction}[theorem]{Construction}

\newcommand{\bin}[2]{\binom{#1}{#2}}
\newcommand{\cG}{\mathcal{G}}
\newcommand{\cC}{\mathcal{C}}
\newcommand{\cB}{\mathcal{B}}
\newcommand{\cR}{\mathcal{R}}
\newcommand{\cD}{\mathcal{D}}
\newcommand{\eps}{\varepsilon}

\title{A combinatorial counting proof of the SOTA lower bound\\
$R(3,k)\ge\bigl(\tfrac12+o(1)\bigr)k^{2}/\log k$}
\author{Derandomization notes\\[0.4em]
\large After Hefty--Horn--King--Pfender, arXiv:2510.19718}
\date{}

\begin{document}
\maketitle

\begin{abstract}
We rewrite the current state-of-the-art lower bound
\[
  R(3,k)\;\ge\;\Bigl(\tfrac12+o(1)\Bigr)\frac{k^{2}}{\log k}
\]
as a pure counting argument over a finite configuration space.
The triangle-cleanup step is deterministic.
No differential equations, no nibble, and no probabilistic language appear:
every estimate is a bound on the number of configurations with a stated combinatorial defect.
This is a \emph{counting derandomization} of~\cite{HHKP2025}, not an explicit construction.
We record separately why a fully explicit derandomization to the same constant remains out of reach.
\end{abstract}

\section{What is being derandomized}

Hefty--Horn--King--Pfender~\cite{HHKP2025} prove
\begin{equation}\label{eq:sota}
  R(3,k)\;\ge\;\Bigl(\tfrac12+o(1)\Bigr)\frac{k^{2}}{\log k}
\end{equation}
by a random ``two bites'' construction: two sparse graphs on $N=n/\log^{2}n$ vertices, their strong product, a random $n$-vertex induced subgraph, and a deterministic greedy deletion of triangles.
Their analysis uses only Chernoff and McDiarmid bounds.

We translate the same argument into counting.
The gain is presentational and conceptual: the proof becomes an existence argument by double counting, matching the classical Erd\H{o}s style, with a fully deterministic cleanup.
The loss is honesty: we do \emph{not} produce an algorithm that outputs the graph, nor an explicit description.

\begin{theorem}[SOTA, counting form]\label{thm:main}
For every $\eps>0$ and all sufficiently large $n$ there exists a triangle-free graph $G$ on $n$ vertices with
\[
  \alpha(G)\;<\;(1+\eps)\sqrt{n\log n}.
\]
Equivalently,~\eqref{eq:sota} holds.
\end{theorem}

\section{Configuration space}

Fix $\eps\in(0,1/100)$ and a large integer $n$.
Set
\begin{align*}
  N &\;=\; \frac{n}{\log^{2} n},
  &
  p &\;=\; \tfrac12\sqrt{\frac{\log n}{n}},
  &
  k &\;=\; (1+\eps)\sqrt{n\log n},
\end{align*}
and fix auxiliary constants $0<\eps_{2}\ll\eps_{1}\ll\eps$.

Let $V_{R},V_{B}$ be disjoint $N$-sets.
Write $\cG_{R}$ (resp.\ $\cG_{B}$) for the set of all graphs on $V_{R}$ (resp.\ $V_{B}$).
Write $\Phi$ for the set of all injections $\varphi\colon[n]\hookrightarrow V_{R}\times V_{B}$.
The \emph{configuration space} is the finite set
\[
  \Omega \;:=\; \cG_{R}\times\cG_{B}\times\Phi.
\]
We equip $\Omega$ with the uniform counting measure: the \emph{fraction} of configurations in a subset $A\subseteq\Omega$ is $|A|/|\Omega|$.

\begin{construction}[Deterministic cleanup]\label{con:cleanup}
Given $(G_{R},G_{B},\varphi)\in\Omega$, form the strong product $G_{1}=G_{R}\boxtimes G_{B}$ on $V_{R}\times V_{B}$:
\[
  (r,b)\sim(r',b')
  \quad\text{iff}\quad
  rr'\in E(G_{R})\ \text{or}\ bb'\in E(G_{B}),
\]
coloured red in the first case and blue in the second.
Let $G_{2}$ be the induced subgraph of $G_{1}$ on $\varphi([n])$.

Build $G=\mathrm{Clean}(G_{R},G_{B},\varphi)$ as follows, in order:
\begin{enumerate}
  \item Inspect red edges of $G_{2}$ in lexicographic order of their $V_{R}$-endpoints; keep an edge iff it creates no monochromatic red triangle with previously kept red edges.
  \item Do the same for blue edges.
  \item From every remaining bichromatic triangle with two blue edges, delete the red edge.
  \item From every remaining bichromatic triangle with two red edges, delete the blue edge.
\end{enumerate}
The output $G$ is triangle-free and isomorphic to a subgraph of $G_{2}$.
\end{construction}

The cleanup depends on no randomness.
Theorem~\ref{thm:main} is equivalent to: the number of configurations whose cleaned graph has an independent set of size $k$ is strictly less than $|\Omega|$.

\section{Binomial tail lemmas (counting form)}

We use only the following elementary bounds on binomial coefficients.
They replace Chernoff/Hoeffding; the proofs are the standard ones via $\bin{m}{t}\le(em/t)^{t}$ and $(1-x)\le e^{-x}$.

\begin{lemma}[Upper tail]\label{lem:tail+}
Let $0\le m\in\mathbb{Z}$, $0<\pi\le 1$, and $X$ a sum of $m$ indicator variables with average value $\mu=m\pi$ under uniform choice of a $\pi$-subset pattern
(equivalently: the fraction of $m$-bit strings of density $\pi$ --- or the fraction of graphs in which a fixed set of $m$ potential edges is present with each edge included in exactly the $\pi$-proportion of graphs --- yielding mean $\mu$).
More concretely, for the uniform distribution on graphs where each edge is present independently in the counting sense that a $\pi$-fraction of graphs contain any fixed edge:
if $Z$ counts the number of edges of a fixed $m$-set that are present, then for every $t\ge 0$,
\[
  \frac{\#\{\text{graphs}: Z\ge\mu+t\}}{\#\{\text{graphs}\}}
  \;\le\;
  \exp\Bigl(-\frac{t^{2}}{2(\mu+t/3)}\Bigr).
\]
\end{lemma}

\begin{lemma}[Lower tail]\label{lem:tail-}
In the same setting, for $0\le t\le\mu$,
\[
  \frac{\#\{\text{graphs}: Z\le\mu-t\}}{\#\{\text{graphs}\}}
  \;\le\;
  \exp\Bigl(-\frac{t^{2}}{2(\mu-t/3)}\Bigr).
\]
\end{lemma}

\begin{lemma}[Bounded-differences count]\label{lem:mcd}
Suppose a function $Z$ on a product space $\prod_{i=1}^{m}\{0,1\}$ changes by at most $c_{i}$ when the $i$-th coordinate flips.
Then for every $t>0$,
\[
  \frac{\#\{Z\ge\mathbb{E}Z+t\}}{2^{m}}
  \;\le\;
  \exp\Bigl(-\frac{t^{2}}{2\sum c_{i}^{2}}\Bigr),
\]
and likewise for the lower tail.
\end{lemma}

\begin{remark}
These are exactly Chernoff and McDiarmid, stated as fraction-of-configurations bounds.
No stochastic process is used.
\end{remark}

\section{Typical configurations}

\begin{definition}[Degree/codegree defect]\label{def:D}
A configuration is in $\cD$ if all of the following hold (with $C=3\sqrt{20}$):
\begin{enumerate}
  \item for every $v\in V_{R}\cup V_{B}$, the fibre size satisfies
    $\bigl||F(v)|-\log^{2}n\bigr|\le\eps_{2}\log^{2}n$,
    where $F(v)$ is the set of selected product-vertices with first (resp.\ second) coordinate $v$;
  \item degrees in $G_{R}$ and $G_{B}$ are $(1\pm\eps_{2})pN$;
  \item codegrees in $G_{R}$ and in $G_{B}$ are at most $C\log n$;
  \item degrees into $G_{2}$ are $(1\pm\eps_{2})pn$;
  \item codegrees into $G_{2}$ are at most $C\log^{3}n$;
  \item projected codegrees across colours are at most $C\log^{3}n$.
\end{enumerate}
\end{definition}

\begin{lemma}[Most configurations are typical]\label{lem:D}
The fraction of configurations outside $\cD$ is $o(1)$.
\end{lemma}

\begin{proof}[Proof sketch]
Each bullet is a degree or codegree count for a fixed vertex or pair.
The mean values are as in~\cite{HHKP2025}: $\mathbb{E}|F(v)|=\log^{2}n$, $\mathbb{E}\deg(v)=pN$, etc.
Lemmas~\ref{lem:tail+}--\ref{lem:tail-} bound the fraction of graphs violating any single constraint by $N^{-\Omega(1)}$.
A union bound over $O(N^{2})$ vertices/pairs leaves an $o(1)$ fraction of bad configurations.
(The passage between a binomial sampling model for $\varphi([n])$ and exact-$n$ injections costs an $N^{O(1)}$ factor absorbed by the same tails; cf.~\cite[Lemma~3.1]{HHKP2025}.)
\end{proof}

\section{Closed pairs inside a large set}

Fix $I\subset[n]$ with $|I|=k$.
For $v\in V_{R}\cup V_{B}$ write $X_{v}(I)=\varphi(I)\cap N_{G_{2}}(v)$.
Partition $V_{R}\cup V_{B}$ by the size of $X_{v}(I)$:
\begin{align*}
  H_{I} &= \{v:|X_{v}(I)|>t_{1}\},
  &
  L_{I} &= \{v:t_{2}<|X_{v}(I)|\le t_{1}\},\\
  M_{I} &= \{v:t_{3}<|X_{v}(I)|\le t_{2}\},
  &
  S_{I} &= \{v:|X_{v}(I)|\le t_{3}\},
\end{align*}
with cutoffs $t_{1}=\sqrt{n\log n}/\log\log n$, $t_{2}=n^{1/4+\eps}$, $t_{3}=n^{2\eps}$.

A pair of vertices in $\varphi(I)$ is \emph{closed} if some common neighbour in $G_{2}$ witnesses it; otherwise \emph{open}.
Open pairs are the only pairs that the cleanup is forced to treat as potential edges of $G$ inside $I$.

\begin{lemma}[Small buckets contribute few closed pairs]\label{lem:buckets}
The fraction of configurations for which there exists an $I$ of size $k$ with
\[
  \sum_{v\in L_{I}}\bin{|X_{v}(I)|}{2}
  +
  \sum_{v\in M_{I}}\bin{|X_{v}(I)|}{2}
  +
  \sum_{v\in S_{I}}\bin{|X_{v}(I)|}{2}
  \;>\;
  3\eps_{1}k^{2}
\]
is $o(1)$.
\end{lemma}

\begin{proof}[Proof sketch]
On $\cD$, the $L_{I}$-sum is $o(k^{2})$ by a deterministic double count: codegrees $\le C\log^{3}n$ force $|L_{I}|$ small, hence $\sum|X_{v}|\le(1+o(1))k$ and each $|X_{v}|\le t_{1}=o(k)$.
For $M_{I}$, an oversized sum produces a bipartite subgraph of $G_{R}\cup G_{B}$ denser than Lemma~\ref{lem:tail+} permits; the number of candidate $(B,I)$ pairs is $\exp(o(n^{3/4-2\eps}))$ while each dense fingerprint has fraction $\exp(-\Omega(n^{3/4-2\eps}))$.
For $S_{I}$, apply Lemma~\ref{lem:mcd} to the bounded-differences function counting closed pairs from small stars; the fraction of bad $I$ is $o(\bin{n}{k}^{-1})$, and summing over $I$ leaves $o(1)$.
\end{proof}

\begin{lemma}[Heavy vertices]\label{lem:heavy}
On $\cD$, for every $I$ of size $k$,
\begin{align*}
  \sum_{v\in H_{I}\cap V_{R}}\bin{|\pi_{B}(X_{v}(I))|}{2}
  &\;\le\;
  \min\Bigl\{\bin{k-|\pi_{R}(I)|}{2},\;
  \bin{pn}{2}+\bin{k-|\pi_{R}(I)|-pn}{2}\Bigr\}+\eps_{1}k^{2},
\end{align*}
and symmetrically for blue.
\end{lemma}

\begin{proof}
Deterministic convexity: the integer vector $(|X_{v}(I)|)_{v\in H_{I}\cap V_{R}}$ is majorized by a vector with entries in $\{0,pn\}$ of the same $\ell_{1}$-mass, up to $\eps_{2}$-errors from~$\cD$.
Projecting to $\pi_{B}$ and using fibre bounds converts the sum of binomials into the displayed minimum.
\end{proof}

\section{Hitting open pairs}

\begin{definition}
For a configuration and a $k$-set $I$, let $\cR_{I}$ be the event that after revealing all edges of $G_{R},G_{B}$ incident to $\pi_{R}(I)\cup\pi_{B}(I)$ from outside, at least
\[
  f(\ell_{R},\ell_{B})-\eps_{1}k^{2}
\]
pairs inside $\pi_{R}(I)$ and $\pi_{B}(I)$ remain unrevealed (open), where $\ell_{R}=|\pi_{R}(I)|$, $\ell_{B}=|\pi_{B}(I)|$, and
\[
  f(\ell_{R},\ell_{B})
  \;=\;
  \bin{\ell_{R}}{2}+\bin{\ell_{B}}{2}
  -\min\Bigl\{\bin{k-\ell_{R}}{2},\bin{pn}{2}+\bin{k-\ell_{R}-pn}{2}\Bigr\}
  -\min\Bigl\{\bin{k-\ell_{B}}{2},\bin{pn}{2}+\bin{k-\ell_{B}-pn}{2}\Bigr\}.
\]
\end{definition}

Lemmas~\ref{lem:D}--\ref{lem:heavy} yield: the fraction of configurations outside $\bigcap_{|I|=k}\cR_{I}$ is $o(1)$.

\begin{lemma}[Open pairs kill independence]\label{lem:kill}
Conditional on $\cR_{I}$ and on $|\pi_{R}(I)|=\ell_{R}$, $|\pi_{B}(I)|=\ell_{B}$, the fraction of extensions of the revealed data for which $I$ is independent in $\mathrm{Clean}(G_{R},G_{B},\varphi)$ is at most
\[
  (1-p)^{f(\ell_{R},\ell_{B})-\eps_{1}k^{2}}
  \;\le\;
  \exp\bigl(-p(f(\ell_{R},\ell_{B})-\eps_{1}k^{2})\bigr).
\]
\end{lemma}

\begin{proof}
Any surviving open red or blue pair, if present as an edge, forces a kept edge of $G$ inside $I$ under the lexicographic cleanup (the same case analysis as~\cite[Lemma~4.1]{HHKP2025}, which is deterministic given the reveal order).
Hence $I$ independent requires all remaining open pairs to be non-edges.
Each such pair is absent from a $(1-p)$-fraction of graphs on that colour, and the pairs are distinct potential edges, so the fraction is at most $(1-p)^{f-\eps_{1}k^{2}}$.
\end{proof}

\begin{lemma}[No large independent set]\label{lem:alpha}
The fraction of configurations for which $\alpha(\mathrm{Clean}(G_{R},G_{B},\varphi))\ge k$ is $o(1)$.
\end{lemma}

\begin{proof}
Write $x_{R}=\ell_{R}/k$, $x_{B}=\ell_{B}/k$.
The fraction of injections with given $(\ell_{R},\ell_{B})$ is at most
\[
  \exp\Bigl(-\frac{2-x_{R}-x_{B}}{2}(1+o(1))\,k\log n\Bigr).
\]
Multiplying by Lemma~\ref{lem:kill} and by $\bin{n}{k}^{-1}$ in the union bound, the same three-case calculus as~\cite[Lemma~4.2]{HHKP2025} with $\beta=\tfrac12$ and $\kappa=1+\eps$ gives, in every case,
\[
  \frac{\#\{\text{configurations with $I$ independent}\}}{|\Omega|}
  \;=\;
  o\Bigl(\bin{n}{k}^{-1}\Bigr).
\]
Summing over at most $\bin{n}{k}$ choices of $I$ leaves an $o(1)$ fraction of bad configurations.
\end{proof}

\section{Proof of Theorem~\ref{thm:main}}

\begin{proof}
By Lemma~\ref{lem:alpha}, a positive fraction of configurations have $\alpha(\mathrm{Clean}(\cdot))<k$.
Any such configuration yields a triangle-free $n$-vertex graph with independence number $<(1+\eps)\sqrt{n\log n}$.
\end{proof}

\section{What this derandomization achieves --- and what it does not}

\begin{center}
\begin{tabular}{lp{0.55\textwidth}}
\hline
Achieved & Not achieved \\
\hline
counting existence proof of~\eqref{eq:sota} & explicit / deterministic poly-time construction \\
deterministic cleanup & removal of binomial tail bounds \\
no DE method, no nibble, no martingales & seed length $o(\sqrt{n}\,(\log n)^{3/2})$ \\
same constant $\tfrac12$ as~\cite{HHKP2025} & matching Alon-style fully explicit graphs at this density \\
\hline
\end{tabular}
\end{center}

\subsection{Barrier to a fully explicit derandomization}

An explicit family of $n$-vertex triangle-free graphs with $\alpha\le(1+\eps)\sqrt{n\log n}$ would be a landmark.
The best explicit constructions remain at $\alpha=O(n^{2/3})$, i.e.\ $R(3,k)=\Omega(k^{3/2})$ (Alon; recent geometric variants).
Spectral methods cannot break this via the Lov\'asz $\vartheta$-function: every triangle-free $n$-vertex graph has $\vartheta=\Omega(\sqrt{n})$, hence any proof that bounds $\alpha$ through $\vartheta$ is stuck at $\Omega(k^{3/2})$.

The counting obstruction is the union bound over $\bin{n}{k}$ potential independent sets: any PRG or limited-independence generator that fools all of them needs seed length
\[
  \Omega\bigl(\log\bin{n}{k}\bigr)\;=\;\Omega\bigl(\sqrt{n}\,(\log n)^{3/2}\bigr),
\]
not $\mathrm{polylog}\,n$.

\subsection{Partial explicit substitutes that fail at constant $\tfrac12$}

\begin{itemize}
  \item Replacing $G_{R},G_{B}$ by polarity graphs of projective planes gives degree $\asymp\sqrt{N}$, the wrong density by a $\sqrt{\log n}$ factor, and independence-number control after the product step is not known to reach $(1+\eps)\sqrt{n\log n}$.
  \item Replacing the random injection $\varphi$ by an explicit extractor / expander walk leaves the two sparse graphs random; the union bound over $k$-sets still requires strong pseudorandomness.
  \item Method of conditional probabilities, applied edge-by-edge, is in principle valid but needs an efficiently computable proxy for the number of surviving bad $k$-sets; no such proxy preserving the constant $\tfrac12$ is known.
\end{itemize}

\section{Conclusion}

The SOTA lower bound admits a clean combinatorial counting proof: average over a finite configuration space, clean triangles deterministically, and bound the number of configurations that leave a $k$-set independent.
This removes probabilistic \emph{language} and all process/DE machinery, while retaining binomial tail estimates as counting lemmas.

A fully explicit derandomization to the same asymptotic constant remains a major open problem, separated from the present argument by the explicit-Ramsey barrier at $\Omega(k^{3/2})$.

\begin{thebibliography}{9}
\bibitem{HHKP2025}
Z.~Hefty, P.~Horn, D.~King and F.~Pfender,
\emph{Improving $R(3,k)$ in just two bites},
arXiv:2510.19718, 2025.

\bibitem{CJMS2025}
M.~Campos, M.~Jenssen, M.~Michelen and J.~Sahasrabudhe,
\emph{A new lower bound for the Ramsey numbers $R(3,k)$},
arXiv:2505.13371, 2025.

\bibitem{Alon1994}
N.~Alon,
\emph{Explicit Ramsey graphs and orthonormal labelings},
Electron.\ J.\ Combin.\ \textbf{1} (1994), \#R12.

\bibitem{Shearer1983}
J.~B.~Shearer,
\emph{A note on the independence number of triangle-free graphs},
Discrete Math.\ \textbf{46} (1983), 83--87.
\end{thebibliography}

\end{document}
